\section{Introducción}
\label{sec:introduccion}
\PARstart{E}{n la era} digital contemporánea, las instituciones de educación superior y centros de investigación dependen de manera fundamental de sus infraestructuras de red de área local (LAN) para articular la totalidad de sus actividades académicas, administrativas e investigativas. Esta dependencia crítica crea una superficie de ataque amplia, compleja y heterogénea, compuesta por cientos o miles de dispositivos finales con niveles de actualización de parches variables, gestionados por equipos de tecnologías de información con recursos de monitoreo y presupuesto limitados. Si bien la conectividad integrada facilita el intercambio eficiente de recursos y la colaboración académica, también expone a la infraestructura institucional a riesgos severos. La presencia de múltiples dispositivos bajo políticas de seguridad no siempre homogéneas genera un entorno altamente vulnerable a la propagación acelerada de amenazas cibernéticas autopropagables, donde la infección de un único host comprometido puede desencadenar efectos en cascada que paralicen el campus completo en cuestión de minutos.

Dentro del catálogo de amenazas informáticas modernas, el ransomware representa una de las amenazas más destructivas para la integridad y disponibilidad de la información organizacional. A diferencia de otros vectores tradicionales de ataque, el ransomware combina el secuestro y cifrado de datos locales con mecanismos dinámicos de propagación lateral capaces de comprometer segmentos completos de red de forma autónoma. Un hito histórico y paradigmático en este ámbito fue el brote masivo del ransomware WannaCry en mayo de 2017, el cual afectó de forma simultánea a más de 230,000 sistemas en 150 países en menos de 24 horas \cite{ghenimi2020wannacry}. Su vector de propagación principal fue la vulnerabilidad EternalBlue (MS17-010), la cual afectaba al protocolo Server Message Block versión 1 (SMBv1) de Windows y permitía la ejecución remota de código sin interacción del usuario final. Las instituciones de educación superior fueron especialmente damnificadas debido a la persistencia de grandes parques de equipos obsoletos, redes planas y sin segmentación lógica entre laboratorios, ausencia de sistemas de monitoreo de tráfico en tiempo real y la exposición irrestricta del puerto TCP 445 entre terminales internas.

Ante la imposibilidad técnica y ética de realizar pruebas de intrusión y propagación real sobre infraestructuras académicas en producción, la simulación basada en agentes (ABM, por sus siglas en inglés) surge como la metodología idónea para estudiar el comportamiento dinámico de estos ataques informáticos. Un modelo de simulación de ciberdefensa permite reproducir escenarios de infección en un entorno virtual controlado y sin riesgo, variando parámetros críticos de protección y contención adaptativa para generar datos cuantitativos que sustenten políticas preventivas. A partir de estas consideraciones, la presente investigación busca responder la siguiente pregunta central de investigación: 

\begin{center}
    \textit{¿En qué medida la combinación de nivel de parche, fuerza del firewall y umbral de contención BDI determina la velocidad de propagación y el alcance final de un ransomware tipo WannaCry en una red LAN universitaria?}
\end{center}

El presente trabajo documenta el desarrollo y evaluación de una simulación espacialmente explícita utilizando GAMA Platform, analizando la efectividad de un protocolo de contingencia activa reactiva.

\subsection{Objetivos de la Investigación}
Para guiar de manera rigurosa el desarrollo del proyecto, se plantearon los siguientes objetivos académicos:

\subsubsection{Objetivo General}
Desarrollar una simulación basada en agentes utilizando la plataforma GAMA Platform que permita analizar la dinámica de propagación de un ransomware con comportamiento de propagación lateral inspirado en WannaCry dentro de una red LAN universitaria, evaluando cuantitativamente la efectividad de una estrategia básica de respuesta automatizada ante una infección masiva.

\subsubsection{Objetivos Específicos}
\begin{itemize}
    \item Diseñar y construir una representación virtual georreferenciada de una red LAN universitaria, mapeando espacialmente los laboratorios y el cableado de comunicaciones mediante herramientas GIS.
    \item Modelar los dispositivos de red (computadoras de escritorio, servidores, switches y firewall) mediante agentes autónomos caracterizados por parámetros individuales de seguridad como niveles de actualización y estado de puertos.
    \item Simular el vector de ingreso de la amenaza externa desde el espacio WAN de Internet atravesando las barreras perimetrales hacia la red interna.
    \item Implementar reglas probabilísticas de transmisión lateral basadas en la exposición del puerto SMB y la efectividad del parcheo de sistemas operativos para replicar la mecánica de propagación de WannaCry.
    \item Evaluar el impacto temporal y cuantitativo de la infección sobre la infraestructura total de equipos conectados.
    \item Diseñar y validar un protocolo de contingencia automatizado basado en aislamiento lógico y blindaje que se active al superar un umbral crítico de infección para limitar la expansión del malware.
\end{itemize}

\subsection{Alcance del Proyecto}
El alcance del presente estudio delimita las fronteras metodológicas de la simulación para garantizar su enfoque académico y su viabilidad técnica:

\begin{itemize}
    \item \textbf{Aspectos Incluidos:} El proyecto contempla el modelado espacial y topológico de una red LAN académica real utilizando datos geográficos reales (Shapefiles), la especificación de un motor probabilístico de propagación de malware fundamentado en la literatura científica de ciberseguridad, y el análisis estadístico del impacto del ransomware con y sin medidas de contingencia activas.
    \item \textbf{Aspectos Excluidos:} No se incluye bajo ninguna circunstancia el desarrollo, almacenamiento o ejecución de código malicioso real, ni la explotación efectiva de vulnerabilidades sobre sistemas físicos o de producción. El proyecto no pretende diseñar un descifrador de archivos ni una herramienta antivirus, sino un modelo computacional educativo y predictivo en un entorno de software controlado.
\end{itemize}

\subsection{Estructura del Documento}
El resto de esta documentación formal se organiza de la siguiente manera: la Sección \ref{sec:estado_arte} describe el estado del arte y la evolución histórica del modelado de propagación de malware en redes complejas, y los fundamentos de la arquitectura cognitiva BDI. La Sección \ref{sec:metodologia} detalla la metodología experimental adoptada, incluyendo la definición de agentes, la integración de datos GIS y el sustento matemático de las ecuaciones de probabilidad de infección. La Sección \ref{sec:arquitectura} expone la arquitectura instanciada de la red LAN universitaria simulada. La Sección \ref{sec:resultados} presenta y discute los resultados cuantitativos de los experimentos de propagación libre y controlada a lo largo de cinco escenarios de defensa. Finalmente, la Sección \ref{sec:conclusiones} resume las principales conclusiones de la investigación y esboza las líneas de trabajo futuro.
